% VI43-DETAILED-CHALLENGE-SOLUTIONS
% VI43-DETAILED-CHALLENGE-01
\subsection*{Challenge VI.43.1}
\begin{solution}
Define \(P\oplus Q\) as the unique \(S\) with \(P+Q\sim O+S\), obtained by the two-line construction.  The Abel--Picard map satisfies \(\alpha_O(P\oplus Q)=\alpha_O(P)+\alpha_O(Q)\).  Degree-one rigidity is exactly what makes \(\alpha_O\) injective: equality of degree-zero point classes implies equality of points.  Therefore
\[
\alpha_O((P\oplus Q)\oplus R)=\alpha_O(P)+\alpha_O(Q)+\alpha_O(R)=\alpha_O(P\oplus(Q\oplus R)),
\]
and injectivity yields associativity.  No coordinate identities are needed.
\end{solution}

% VI43-DETAILED-CHALLENGE-02
\subsection*{Challenge VI.43.2}
\begin{solution}
For an arbitrary origin \(O\), the operation is characterized by \(P+Q\sim O+(P\oplus_O Q)\).  If \(F\) is a flex and \(\oplus_F\) is the flex-origin law, then the Abel--Picard identifications differ by translation of divisor classes:
\[
[P-O]=[P-F]+[F-O].
\]
Thus the two group structures are conjugate by a translation on the underlying genus-one curve; the identity moves from \(F\) to \(O\).  The one-line rule \(P+Q+R\sim3F\) uses the special equality \(H\sim3F\), unavailable for general \(O\), which is why the arbitrary-origin construction needs two lines.
\end{solution}

% VI43-DETAILED-CHALLENGE-03
\subsection*{Challenge VI.43.3}
\begin{solution}
Choose projective coordinates sending the flex to \([0:1:0]\) and its tangent to \(z=0\).  The flex condition forces the intersection with \(z=0\) to be triple at that point, eliminating the unwanted cubic terms along the line at infinity.  An affine change then completes the square in \(y\) (characteristic not \(2\)) and shifts/scales \(x\) to remove the quadratic term (characteristic not \(3\)).  After rescaling one obtains
\[
y^2z=x^3+Axz^2+Bz^3.
\]
Throughout, the chosen flex remains the point at infinity and \(z=0\) remains its tangent, so it is the identity for the standard chord law.
\end{solution}

% VI43-DETAILED-CHALLENGE-04
\subsection*{Challenge VI.43.4}
\begin{solution}
Take a Weierstrass family with discriminant tending to zero.  At a simple discriminant zero the limiting cubic is nodal; normalization produces two preimages of the node and the smooth locus becomes \(\mathbb P^1\setminus\{0,\infty\}\simeq\mathbb G_m\).  At a cuspidal degeneration the two branches coalesce, leaving one normalization point and smooth locus \(\mathbb A^1\simeq\mathbb G_a\).  Tracking the secant parameter shows multiplication in the nodal case and addition in the cuspidal case.  The discriminant therefore marks the boundary where the elliptic group degenerates to multiplicative or additive type.
\end{solution}

% VI43-DETAILED-CHALLENGE-05
\subsection*{Challenge VI.43.5}
\begin{solution}
For the standard quadratic transformation, a degree-\(3\) cubic with multiplicities \(m_i\) at the three base points transforms to degree \(6-\sum m_i\).  If all three base points lie simply on a smooth cubic, the transform is again degree \(3\) and birational to the original genus-one curve; if fewer lie on it, the plane degree rises and singularities must account for the unchanged normalization genus.  The divisor-class shadow is \(H\mapsto2H-E_1-E_2-E_3\).  Thus the intrinsic Picard/group-law data survive on the normalization while the plane embedding and hyperplane class change birationally.
\end{solution}

